% GENERATED FILE. Edit canonical lessons, then run npm run generate:books.
% canonical-source-hash: fnv1a64:be4a2127937867e3
% canonical-lessons: HI-C458-tanav, HI-C458-fikr, HI-C458-bojh, HI-C458-tarapna, HI-C458-asvasan

\chapter{Worry and Relief}
\label{ch:hi-worry-and-relief}
\hlchaptermodality{458}

\begin{chapteropening}
\textbf{By the end of this chapter:} \emph{I can say what is worrying someone and when it eased.}

Complete the last lesson of chapter 458: I can say what is worrying someone and when it eased.
\end{chapteropening}

\section[\texorpdfstring{\hi{तनाव} (tanāv)}{तनाव (tanāv)}]{\hi{तनाव} (tanāv) — stress; tension}

\label{lesson:HI-C458-tanav}



\begin{quote}
Before the new one: say the Hindi for cowardly, then the Hindi for to shudder.
\end{quote}

\subsection*{You'll want to know: \hi{तनाव}}
\textbf{\hi{तनाव}} — \emph{tanāv} — \textquotedblleft{}stress; tension\textquotedblright{}.

\begin{quote}
\hi{नौकरी} \hi{के} \hi{तनाव} \hi{से} \hi{अरुण} \hi{की} \hi{नींद} \hi{चली} \hi{गई।}
\end{quote}

\begin{quote}
\emph{naukrī ke tanāv se Aruṇ kī nīnd calī gaī.}
\end{quote}

\begin{quote}
The stress of his job has cost Arun his sleep.
\end{quote}

Masculine.

\begin{cousinweb}
\textbf{\hi{तनाव}} (\emph{tanāv}) is probably a noun built on the root of \emph{tānnā}, \textquotedblleft{}to stretch, pull, tighten\textquotedblright{}, which CDIAL lists under Sanskrit \emph{tānáyati} (CDIAL 5762): a rope pulled tight, and then a mind. CDIAL gives no row for \emph{tanāv} itself.
\end{cousinweb}

\begin{grammarlens}[title={What you've built}]
One more word for describing what you lived and felt.
\end{grammarlens}

\subsection*{Your turn}
\begin{itemize}
\raggedright
  \item \emph{Say it:} \emph{tanāv}
  \item \emph{Say it:} \emph{tanāv}, once more
  \item \emph{Say it:} \emph{naukrī ke tanāv se Aruṇ kī nīnd calī gaī.} — The stress of his job has cost Arun his sleep.
  \item \emph{Recall:} say the Hindi for cowardly, then the Hindi for to shudder, then say \emph{tanāv} again
\end{itemize}

\begin{tcolorbox}[breakable,colback=teal!4,colframe=teal!35!black,title={Before you move on}]
What does \hi{तनाव} mean? (\textquotedblleft{}Stress; tension\textquotedblright{}.) Say it once more.
\end{tcolorbox}
\section[\texorpdfstring{\hi{फ़िक्र} (fikr)}{फ़िक्र (fikr)}]{\hi{फ़िक्र} (fikr) — concern; a worry}

\label{lesson:HI-C458-fikr}



\begin{quote}
Before the new one: say the Hindi for to shudder, then the Hindi for stress.
\end{quote}

\subsection*{You'll want to know: \hi{फ़िक्र}}
\textbf{\hi{फ़िक्र}} — \emph{fikr} — \textquotedblleft{}concern; a worry\textquotedblright{}.

\begin{quote}
\hi{माँ} \hi{को} \hi{हमेशा} \hi{बच्चों} \hi{की} \hi{फ़िक्र} \hi{रहती} \hi{है।}
\end{quote}

\begin{quote}
\emph{māṁ ko hameśā baccoṁ kī fikr rahtī hai.}
\end{quote}

\begin{quote}
Mother always worries about the children.
\end{quote}

Feminine.

\begin{cousinweb}
\textbf{\hi{फ़िक्र}} (\emph{fikr}) is usually traced to Arabic \emph{fikr}, \textquotedblleft{}thought, reflection\textquotedblright{}, which reached Hindi through Persian; in Hindi the thought became a worry.
\end{cousinweb}

\begin{grammarlens}[title={What you've built}]
One more word for describing what you lived and felt.
\end{grammarlens}

\subsection*{Your turn}
\begin{itemize}
\raggedright
  \item \emph{Say it:} \emph{fikr}
  \item \emph{Say it:} \emph{fikr}, once more
  \item \emph{Say it:} \emph{māṁ ko hameśā baccoṁ kī fikr rahtī hai.} — Mother always worries about the children.
  \item \emph{Recall:} say the Hindi for to shudder, then the Hindi for stress, then say \emph{fikr} again
\end{itemize}

\begin{tcolorbox}[breakable,colback=teal!4,colframe=teal!35!black,title={Before you move on}]
What does \hi{फ़िक्र} mean? (\textquotedblleft{}Concern; a worry\textquotedblright{}.) Say it once more.
\end{tcolorbox}
\section[\texorpdfstring{\hi{बोझ} (bojh)}{बोझ (bojh)}]{\hi{बोझ} (bojh) — a burden; a weight on the mind}

\label{lesson:HI-C458-bojh}



\begin{quote}
Before the new one: say the Hindi for stress, then the Hindi for concern.
\end{quote}

\subsection*{You'll want to know: \hi{बोझ}}
\textbf{\hi{बोझ}} — \emph{bojh} — \textquotedblleft{}a burden; a weight on the mind\textquotedblright{}.

\begin{quote}
\hi{सच} \hi{कहने} \hi{के} \hi{बाद} \hi{अरुण} \hi{के} \hi{मन} \hi{से} \hi{बोझ} \hi{उतर} \hi{गया।}
\end{quote}

\begin{quote}
\emph{sac kahne ke bād Aruṇ ke man se bojh utar gayā.}
\end{quote}

\begin{quote}
Once he had told the truth, a weight lifted from Arun's mind.
\end{quote}

Masculine.

\begin{cousinweb}
\textbf{\hi{बोझ}} (\emph{bojh}) is a tadbhava: CDIAL lists Hindi \emph{bojh}, \textquotedblleft{}load\textquotedblright{}, under Sanskrit \emph{vahyá} (CDIAL 11465). A load on the back and a load on the mind are both \emph{bojh}.
\end{cousinweb}

\begin{grammarlens}[title={What you've built}]
One more word for describing what you lived and felt.
\end{grammarlens}

\subsection*{Your turn}
\begin{itemize}
\raggedright
  \item \emph{Say it:} \emph{bojh}
  \item \emph{Say it:} \emph{bojh}, once more
  \item \emph{Say it:} \emph{sac kahne ke bād Aruṇ ke man se bojh utar gayā.} — Once he had told the truth, a weight lifted from Arun's mind.
  \item \emph{Recall:} say the Hindi for stress, then the Hindi for concern, then say \emph{bojh} again
\end{itemize}

\begin{tcolorbox}[breakable,colback=teal!4,colframe=teal!35!black,title={Before you move on}]
What does \hi{बोझ} mean? (\textquotedblleft{}A burden; a weight on the mind\textquotedblright{}.) Say it once more.
\end{tcolorbox}
\section[\texorpdfstring{\hi{तड़पना} (taṛapnā)}{तड़पना (taṛapnā)}]{\hi{तड़पना} (taṛapnā) — to writhe; to be in anguish}

\label{lesson:HI-C458-tarapna}



\begin{quote}
Before the new one: say the Hindi for concern, then the Hindi for a burden.
\end{quote}

\subsection*{You'll want to know: \hi{तड़पना}}
\textbf{\hi{तड़पना}} — \emph{taṛapnā} — \textquotedblleft{}to writhe; to be in anguish\textquotedblright{}.

\begin{quote}
\hi{दाँत} \hi{के} \hi{दर्द} \hi{से} \hi{अरुण} \hi{पूरी} \hi{रात} \hi{तड़पा।}
\end{quote}

\begin{quote}
\emph{dā̃t ke dard se Aruṇ pūrī rāt taṛpā.}
\end{quote}

\begin{quote}
Arun writhed all night with toothache.
\end{quote}

Past (he …): \hi{वह} \hi{दर्द} \hi{से} \hi{तड़पा}, \emph{vah dard se taṛpā}.

\begin{cousinweb}
\textbf{\hi{तड़पना}} (\emph{taṛapnā}) is a tadbhava: CDIAL lists Hindi \emph{taṛapnā}, \textquotedblleft{}to toss, roll about, be anxious\textquotedblright{}, under \emph{*taḍapphaḍ}, a form it reconstructs (CDIAL 5634).
\end{cousinweb}

\begin{grammarlens}[title={What you've built}]
One more word for describing what you lived and felt.
\end{grammarlens}

\subsection*{Your turn}
\begin{itemize}
\raggedright
  \item \emph{Say it:} \emph{taṛapnā}
  \item \emph{Say it:} \emph{taṛapnā}, once more
  \item \emph{Say it:} \emph{dā̃t ke dard se Aruṇ pūrī rāt taṛpā.} — Arun writhed all night with toothache.
  \item \emph{Recall:} say the Hindi for concern, then the Hindi for a burden, then say \emph{taṛapnā} again
\end{itemize}

\begin{tcolorbox}[breakable,colback=teal!4,colframe=teal!35!black,title={Before you move on}]
What does \hi{तड़पना} mean? (\textquotedblleft{}To writhe; to be in anguish\textquotedblright{}.) Say it once more.
\end{tcolorbox}
\section[\texorpdfstring{\hi{आश्वासन} (āśvāsan)}{आश्वासन (āśvāsan)}]{\hi{आश्वासन} (āśvāsan) — an assurance; comforting words}

\label{lesson:HI-C458-asvasan}



\begin{quote}
Before the new one: say the Hindi for a burden, then the Hindi for to writhe.
\end{quote}

\subsection*{You'll want to know: \hi{आश्वासन}}
\textbf{\hi{आश्वासन}} — \emph{āśvāsan} — \textquotedblleft{}an assurance; comforting words\textquotedblright{}.

\begin{quote}
\hi{डाक्टर} \hi{के} \hi{आश्वासन} \hi{से} \hi{मीरा} \hi{की} \hi{चिंता} \hi{कम} \hi{हुई।}
\end{quote}

\begin{quote}
\emph{ḍākṭar ke āśvāsan se Mīrā kī cintā kam huī.}
\end{quote}

\begin{quote}
The doctor's assurance eased Mira's worry.
\end{quote}

Masculine.

\begin{cousinweb}
\textbf{\hi{आश्वासन}} (\emph{āśvāsan}) is a tatsama, borrowed from Sanskrit \emph{āśvāsana}, \textquotedblleft{}reviving; consoling, encouraging\textquotedblright{} (Monier-Williams).
\end{cousinweb}

\begin{grammarlens}[title={What you've built}]
One more word for describing what you lived and felt.
\end{grammarlens}

\subsection*{Your turn}
\begin{itemize}
\raggedright
  \item \emph{Say it:} \emph{āśvāsan}
  \item \emph{Say it:} \emph{āśvāsan}, once more
  \item \emph{Say it:} \emph{ḍākṭar ke āśvāsan se Mīrā kī cintā kam huī.} — The doctor's assurance eased Mira's worry.
  \item \emph{Recall:} say the Hindi for a burden, then the Hindi for to writhe, then say \emph{āśvāsan} again
\end{itemize}

\begin{tcolorbox}[breakable,colback=teal!4,colframe=teal!35!black,title={Before you move on}]
What does \hi{आश्वासन} mean? (\textquotedblleft{}An assurance; comforting words\textquotedblright{}.) Say it once more.
\end{tcolorbox}
